\documentclass{beamer}

\usepackage{amsmath}
\usepackage{amsfonts}
\usepackage{amssymb}
\usepackage{textpos}
\usepackage[orientation=portrait, size=a0,scale=1.21]{beamerposter}
\usepackage{ragged2e}
\usepackage{graphicx}

\usepackage{subfigure}


\usetheme{Berlin}
%\usecolortheme{dolphin}

\setbeamertemplate{headline}{}
\setbeamertemplate{footline}{}


\usepackage{xcolor}

% Redefine block environment background to white
\setbeamercolor{block body}{bg=white}
\setbeamercolor{block title}{bg=purple}
\setbeamercolor{title}{bg=purple}


\begin{document}
%\begin{textblock}{100}(0.1,1.6)

\title[]{\Huge Container-based high-performance computing for modeling large-scale biological neural networks  \\[-0.1 in]}
\author[]{\LARGE \textbf{Oliver James} and \textbf{Sungho Hong} \\[-0.1 in] }
\institute[]{\Large Center for Cognition and Sociality, Institute for Basic Science}
\date{}



\begin{frame}
\vspace{-0.3 in}
\maketitle

 \begin{columns}
        % First column
        \begin{column}{0.49\linewidth}
        \vspace{0 in}
            \begin{block}{\LARGE Introduction}
            \justifying
        {\large 1. Modeling large-scale, biophysically detailed neural networks increasingly relies on high-performance computing (HPC), but deploying such simulations across diverse supercomputing environments remains challenging due to software incompatibilities and complex dependencies \\[0.3 cm]
2. Singularity containers address these issues by encapsulating the entire simulation environment, enabling neuroscientists to achieve portability and reproducibility without requiring root privileges on shared HPC systems.\\[0.3 cm]
3. Despite these advantages, container adoption in computational neuroscience is still limited, and practical demonstrations—such as our cerebellar mossy fiber–granule cell network simulations—are needed to showcase their effectiveness for scalable, reproducible research.
               }
            \end{block}
            \vspace{0.5cm}
            \begin{block}{\LARGE Containerization for Scientific Computing }
            \centering
            \vspace{0.5cm}
             %\includegraphics[scale=2.2 ]{/home/olive/Desktop/figs/tex/Paradigm_ReportOnly.eps}
            \end{block}

            \begin{block}{\LARGE Our Computational Framework (N=27)}
            \centering
               %\includegraphics[scale=1.7 ,angle=-90]{/home/olive/Desktop/figs/tex/Poster_behavior.eps}
            \end{block}

            A slight change in mask contrast induces blindsight-like behavioral performance.

             \vspace{0.5cm}
            \begin{block}{\LARGE Case Study: Cerebellar Mossy Fiber–Granule Cell Network}
            \centering
            \vspace{1 cm}
            %\includegraphics[scale=1.4 ,angle=90]{/home/olive/Desktop/figs/tex/Decoder.eps}

\vspace{3 cm}

             %\includegraphics[width=0.99\textwidth]{/home/olive/Desktop/figs/tex/PosterEEG.eps}
            \end{block}

        There exists a \emph{switch} in the spatio-temporal dynamics of stimulus representation, indicating the involvement of distinct neural mechanisms.  \\[0.3 cm]

       %In visible conditions, the transition from dynamic to stable coding occurs at later time points, whereas in invisible conditions, this transition manifests earlier.



          %Visible conditions are characterized by strong dynamic neural coding, whereas invisible conditions exhibit early stable neural representations of the stimulus.

        \end{column}

%################### Second column ####################
        \begin{column}{0.49\textwidth}
            \vspace{-0.5 cm}
            \centering
            %\includegraphics[scale=1.74,angle=-90]{/home/olive/Desktop/figs/tex/Fig2c.eps}

       \vspace{0.5 cm}

        \justifying
        In visible conditions, robust dynamic neural representations in sensory and central areas suggest that this mechanisms may be sufficient to support initial stages of conscious awareness.


        During invisible conditions, with weak dynamic coding, early stable neural representations of the stimulus emerge from the frontal cortex, implicating a potential role of frontal recurrent neural circuit for conscious awareness.

       \vspace{0.5 cm}

            \begin{block}{\LARGE Performance Across HPC Environments}
            \centering
            %\includegraphics[width=0.92\textwidth]{/home/olive/Desktop/figs/tex/Fig3.eps}

            \justifying
            Source space analysis reveals the initiation of stable neural representations, progressing from frontal to visual cortices, during conditions of perceptual invisibility.


            \end{block}

            \begin{block}{\LARGE Reproducibility & Portability}

            \vspace{0.5 cm}

           A five-cell mesoscopic canonical model is proposed for each region within the source space

             \vspace{0.5 cm}

            \centering
                 \includegraphics[scale=1.7]{/home/olive/Desktop/nature_neuroscience/Entire.eps}




            \end{block}

\begin{figure}[htbp]
    \centering
    \subfigure[$ \gamma_{\text{PYR}}$]{%
        \includegraphics[width=0.27\textwidth]{/home/olive/Desktop/figs/tex/Poster_pyr.eps}

    } \quad
    \subfigure[$ \gamma_{\text{VIP+}}$]{%
        \includegraphics[width=0.27\textwidth]{/home/olive/Desktop/figs/tex/Poster_vip.eps}

    }\quad
    \subfigure[$ \gamma_{\text{PV+}}$]{%
        \includegraphics[width=0.31\textwidth]{/home/olive/Desktop/figs/tex/Poster_pv_colorbar.eps}
 }
\end{figure}

             \begin{block}{\normalsize Acknowledgment}
             \justifying

This work was supported by IBS-R001-D2-2023-a00

               \end{block}

        \end{column}
    \end{columns}

\end{frame}
%\end{textblock}

\end{document}
